\clearpage
\section*{Appendix Y. Quarter-turn symmetry and certified signed pressure}
\addcontentsline{toc}{section}{Appendix Y. Quarter-turn symmetry and signed pressure}
\noindent\textbf{Scope.} This appendix gives a source-independent exact periodic-pressure test of the quarter-turn symmetry used by H8 localization and a fourfold angular-mode error bound. It does not authenticate the Case~006 H8 data or construct a nonlinear unforced solution.

\paragraph{Proposition Y.1 (Fourfold selection with a symmetry caveat).}
On the square periodic three-torus let $Q(x,y,z)=(-y,x,z)$. If a smooth divergence-free velocity satisfies $u(Qx)=Qu(x)$, its unique mean-zero Poisson pressure
$\pi[u]=(-\Delta)^{-1}_0\partial_i\partial_j(u_i u_j)$ is $Q$-invariant. Hence its circular trace has Fourier coefficients $\pi_m=0$ unless $m\in4\mathbb Z$. The quarter-turn average of a vector field has this covariance. However, if only the prepared carrier $U$ is equivariant, the conclusion is valid for $\pi[U]$ but does \emph{not} automatically hold for $\pi[U+w]$: the exact polarization
$\pi[U+w]=\pi[U]+2\mathcal Q(U,w)+\pi[w]$ can break the symmetry when $w$ does. Averaging a proposed nonlinear correction does not generally preserve the correction equation.
\begin{proof}
The scalar Poisson source is invariant under $Q$ and the periodic mean-zero inverse is unique. Quarter-turn invariance of the scalar trace forces its Fourier support to lie in $4\mathbb Z$. The last distinction follows from pressure polarization.
\end{proof}

\paragraph{Theorem Y.2 (Fourfold-symmetric pressure work with opposite signs).}
Let
\[
\psi_\alpha(x,y)=\cos x+\cos y+
\alpha\sin x\sin y(\cos x-\cos y),\qquad
u_\alpha=(\partial_y\psi_\alpha,-\partial_x\psi_\alpha,0)
\]
on $(\mathbb R/2\pi\mathbb Z)^3$. Each $u_\alpha$ is smooth, divergence-free and $C_4$-equivariant. Reconstruct its \emph{actual} mean-zero periodic Poisson pressure $\pi_\alpha$ and consider the embedded unit cylindrical region
$\Omega=\{x^2+y^2<1\}\times\mathbb T_z$. Define work \emph{per unit axial length},
\[
\mathcal W(\alpha)=\frac1{2\pi}
\int_{\partial\Omega}\pi_\alpha u_\alpha\cdot n\,dS
=\int_0^{2\pi}\pi_\alpha(\cos\theta,\sin\theta)\,
u_\alpha\cdot e_R\,d\theta.
\]
For $\alpha=3/10$ there is the exact finite Bessel formula
\begin{equation}\label{eq:Y-bessel}
\boxed{\mathcal W(3/10)=2\pi\sum_{n\in\mathcal N}c_n\frac{J_1(\sqrt n)}{\sqrt n},\qquad
\mathcal N=\{1,5,9,13,17,29,37,41\},}
\end{equation}
with
\[
\begin{array}{c|rrrrrrrr}
n&1&5&9&13&17&29&37&41\\ \hline
c_n&-2427/32000&21/100&-10791/64000&3/640&
1047/32000&-81/16000&27/8000&-81/64000
\end{array}
\]
and the \emph{rigorously rationally enclosed} signs
\begin{equation}\label{eq:Y-sign}
\boxed{-0.0079871<\mathcal W(3/10)<-0.0079870,\qquad
\mathcal W(-3/10)=-\mathcal W(3/10)>0.}
\end{equation}
The full three-dimensional work equals $2\pi\mathcal W(\alpha)$.
Neither snapshot is claimed to solve the time-dependent unforced Navier--Stokes equation.
\begin{proof}
Since
$\sin x\sin y(\cos x-\cos y)=
(\sin2x\sin y-\sin x\sin2y)/2$,
the stream function is invariant under $(x,y)\mapsto(-y,x)$.
Writing $\psi_\alpha=\sum_a s_a e^{ia\cdot(x,y)}$ gives rational finite Fourier coefficients $s_a$ and velocity coefficients $i(a_y,-a_x)s_a$. For $k\ne0$, the complete periodic Poisson pressure has the finite Fourier coefficient
\[
\widehat\pi_\alpha(k)
=\frac{1}{|k|^2}\sum_{a+b=k}
(k_xa_y-k_ya_x)(k_xb_y-k_yb_x)s_as_b,\qquad
\widehat\pi_\alpha(0)=0.
\]
On the unit circle use
$\int_0^{2\pi}e^{iq\cdot e_R}e_R\,d\theta
=2\pi iJ_1(|q|)\,q/|q|$.
Exact rational convolution and grouping by $n=|q|^2$ produce the eight coefficients displayed.
The expansion
\[
\frac{J_1(\sqrt n)}{\sqrt n}
=\frac12\sum_{j=0}^\infty\frac{(-n/4)^j}{j!(j+1)!}
\]
has total truncation error after $j=20$ at most
\[
E_{20}=\frac12\sum_{n\in\mathcal N}|c_n|
\frac{(n/4)^{21}}{21!\,22!}<3\times10^{-23};
\]
the alternating tails are decreasing there because $n\le41$.
Using rational bounds $314159/100000<\pi<314160/100000$ establishes the strict interval \eqref{eq:Y-sign} with no floating-point sign assumptions.
Finally $\psi_{-\alpha}(x,y)=\psi_\alpha(y,x)$:
coordinate exchange reverses the oriented stream velocity while retaining the corresponding Poisson pressure, and consequently reverses $\mathcal W$.
\end{proof}

\paragraph{Proposition Y.3 (Quantitative $m=4$ dominance gate).}
For a circular reference interface $R=a(\eta)$ under Appendix~X's leading lateral solenoidality, let $q(\theta,\eta)=V_\perp\cdot e_R$ have $q_0(\eta)=0$. If the \emph{full} leading pressure trace $\Pi$ is $C_4$-invariant, then
\[
C_\pi=4\pi\int a(\eta)\operatorname{Re}
(\Pi_4(\eta)\overline{q_4(\eta)})\,d\eta+\mathcal E_{\ge8}
\]
and
\begin{equation}\label{eq:Y-angular-tail}
\boxed{|\mathcal E_{\ge8}|\le\frac1{64}
\int a(\eta)\|\partial_\theta^2\Pi(\cdot,\eta)\|_{L^2_\theta}
\|q(\cdot,\eta)\|_{L^2_\theta}\,d\eta.}
\end{equation}
The bound also applies to carrier self-pressure when only $U$ is $C_4$-equivariant. It cannot be applied to \emph{full} pressure unless the complete velocity and nonlinear correction retain the symmetry.
\begin{proof}
Only angular frequencies $m\in4\mathbb Z$ occur in $\Pi$, and $q_0=0$. After the $m=\pm4$ modes, $|m|\ge8$, so $|m|^{-2}\le1/64$; Parseval and Cauchy--Schwarz give \eqref{eq:Y-angular-tail}.
\end{proof}

\paragraph{Conditional sign test and remaining frontier.}
Let $C_4$ denote the explicit $m=\pm4$ term and $B_8$ the bound on the right of \eqref{eq:Y-angular-tail}. If a single authenticated H8 solution supplies $C_{\rm other}\le C_{\rm other}^+$ and
\[
C_4+B_8+C_{\rm other}^+\le -M_0(R_0-p/2),
\]
then the \emph{necessary inward-energy coefficient gate} of Appendix W is certified. Conversely, a lower bound $C_4-B_8+C_{\rm other}^->-M_0(R_0-p/2)$ excludes that source/solution under the stated hypotheses. Neither sign test proves that a nonlinear correction exists. The original H8 vector potential, cutoff, shapers and overlaps, frame, density, stress centre, complete pressure, exterior correction and terminal retention remain unauthenticated together. The full unforced Navier--Stokes finite-time blow-up theorem remains \textbf{OPEN}. The Fourier selection and Bessel techniques are classical, not new general theorems.
